% page 341 of the facsimile (image p-340 of the scan)
% block 165.1 x 267.9 mm ; left margin 36.9 mm
\pgc{15.240mm}{0.000mm}{
\l{}
\l{}
\l{}
\l{}
\l{\cel{51}333.}
\l{}
\l{leda. Qua aspektas desagreabla pro ne-beleso. Qua desplezas}
\l{\cel{3}etikale o morale. - FI.}
\l{}
\l{ledro. Pelo dika, deprenita de la karno di ula animali (bovulo,}
\l{\cel{3}bovino, bovyuno, bufalo, e c.), senpiligita per tanago, e}
\l{\cel{3}preparita por uzi diversa. - DE.}
\l{}
\l{\sou{lego}. I. Ago-regulo impozata da autoritatozo supera. - II. Ago-}
\l{\cel{3}regulo impozata a la homo da lua raciono, lua koncienco. -}
\l{\cel{3}III. (+\sou{leyo}) Regulo konstanta, universala, a qua omna feno-}
\l{\cel{3}meni di naturo esas obedienda. - Rezultajo de analizar exakte}
\l{\cel{3}e precize la cirkonstanci qui esas elementi en la produkto}
\l{\cel{3}di la fenomeni, edinterligar ta cirkonstanci per la relati}
\l{\cel{3}normala di inter-sucedo e di inter-simileso. (Auguste Comte).}
\l{\cel{3}DEFIRS.}
\l{}
\l{\sou{legacar}. (\sou{trans.,}\cel{1}ad). Donacar (ad ulu) per testamento qua atri-}
\l{\cel{3}buas ad un od a plura personi parto di sua havaji o nur parto}
\l{\cel{3}de olci. - DEFIS.}
\l{}
\l{\sou{legato}. I. Delegito di la imperiestri Romana, qua havis kom tasko}
\l{\cel{3}reprezentar ici en la provinci. - II. (\sou{Eklezio katol}.) Kar-}
\l{\cel{3}dinalo delegita da papo por guvernar un ek la provinci di la}
\l{\cel{3}eklezio katolika. - DEFIRS.}
\l{}
\l{\sou{legendo}. Serio de rakonti populala, qui havas maxim-multa-kaze}
\l{\cel{3}fundo exakta, ma developita, transformita da la tradicioni.}
\l{\cel{3}- DEFIRS.}
\l{}
\l{\sou{legiono}. I. (l)(\sou{epoki antiqua, en Roma}). Armeo-korpo qua konsis-}
\l{\cel{3}tis ye infantrio e kavalrio. (2) En Francia, lor la rejo}
\l{\cel{3}Francisko I. (cirkum la yaro l5l5) : armeo-korpo permananta.}
\l{\cel{3}- II. (\sou{metaf}.)Grupo de individui qui esforcas atingar la}
\l{\cel{3}sama skopo. - III. (en\cel{1}\sou{Francia}): Legiono di Honoro : merit-}
\l{\cel{3}ordeno, civila e militala, institucita da Bonaparte, konsulo}
\l{\cel{3}rang-unesma. - DEFIRS.}
\l{}
\l{\sou{legitima}. Konsakrita da la lego.}
\l{}
\l{\sou{legumo.}\cel{2}Parto di planto koliita por uzesar kom nutrivo : lensi,}
\l{\cel{3}fazeoli, pizi, asparagi, e c. DEFIS.}
\l{}
\l{\sou{leguminoso}. (\sou{bot.}) Di qua la frukto esas shelo qua kontenas}
\l{\cel{3}maxim-multa-kaze legumo (pizo, fabo, fazeolo). - DEFIS.}
\l{}
\l{\sou{lejera}.\cel{2}- Qua pezas nur poke. - FIS.}
\l{}
\l{\sou{lekar}.\cel{2}(\sou{trans.)}\cel{3}Pasigar sua lango sur (ulo). - DEFI.}
\l{}
\l{\sou{lektar}.\cel{2}(\sou{trans.})\cel{2}Dicarnar, en imprimuro o skriburo, la son-uni}
\l{\cel{3}reprezentita per literi. - II. Konoceskar la kontenajo di}
\l{\cel{3}skriburo, di libro. - III. Konocigar, savigar,da altra per-}
\l{\cel{3}soni la kontenajo di skriburo, di libro, per pronuncar koram}
\l{\cel{3}li to quo esas skribita, imprimita. -IV. (metaf.) Divinar la}
\l{\cel{3}penso da ulu, lua sentimenti, ek lua konduto, fiziognomio,}
\l{\cel{3}aspekto extera, e c. FIRS.}
}
